\documentclass[10pt,a4paper]{amsart}
\usepackage[margin=19mm]{geometry}
\usepackage{amsmath,amssymb,mathtools}
\usepackage[scr=rsfs,cal=euler]{mathalfa}
\usepackage{xfrac}
\usepackage[unicode,hidelinks,bookmarks=false]{hyperref}
\newcommand{\ie}{i.e.\ }
\newcommand{\cf}{cf.\ }
\newcommand{\resp}{respectively\ }
\newcommand{\Subsection}{\subsection}
\providecommand{\nobreakdash}{\nobreak\hbox{-}}
\setlength{\emergencystretch}{2em}
\multlinegap0pt
\usepackage{amssymb}
\usepackage{mathtools}
\usepackage{xcolor}
\newtheorem{lemma}{Lemma}
\newtheorem{proposition}{Proposition}
\newcommand{\FF}{\mathcal{F}}
\newcommand{\LL}{\mathcal{L}}
\newcommand{\OO}{\mathcal{O}}
\renewcommand{\P}{\mathbb{P}}
\newcommand{\PP}{\mathcal{P}}
\newcommand{\gDD}{\operatorname{gDD}}
\newcommand{\rank}{\operatorname{rank}}
\title{Osculating Geometry and Higher-Order Distance Loci}
\author{Sandra Di Rocco, Kemal Rose, Luca Sodomaco}
\date{Source: arXiv:2507.02823v4 (CC BY 4.0)}
\begin{document}
\maketitle

\noindent\emph{Source and licence.} Osculating Geometry and Higher-Order Distance Loci. Version arXiv:2507.02823v4 (CC BY 4.0), distributed by arXiv under the Creative Commons Attribution 4.0 International licence, http://creativecommons.org/licenses/by/4.0/, as recorded in the arXiv licence metadata for this exact version and shown on the abstract page https://arxiv.org/abs/2507.02823v4. Full licence notice: \texttt{licenses/arxiv-2507.02823-CC-BY-4.0.txt}.\par\smallskip

\noindent\textit{Editorial scope.} Counted unit: the article's proposition surface (source0.tex lines 1521--1533). The article's complete original proof of it is delivered with it (source0.tex lines 1534--1579, one \texttt{proof} environment of 46 lines). No mathematics is written, repaired or re-derived: the statement and the proof are the article's own lines, in the article's own notation, and no retained line is substituted or re-flowed.  Retained uncounted as the source-local context the proof needs: lines 1442--1447; lines 1450--1452; lines 1468--1473; lines 1479--1491; lines 1501--1503.  Nothing was omitted from any retained span, so the package carries no omission record.  No retained line contains a citation, so the wrapper carries no bibliography and no bibliography entry.  No retained line contains a figure environment, an image-inclusion command or drawing code, so no image is delivered and the package declares no asset dependency.  Editorial formatting only, in addition: an AMS \texttt{amsart} preamble that restates the article's own declarations and macros used by the retained lines, namely the environments \texttt{lemma}, \texttt{proposition} on separate counters, together with the standard packages those lines need.  The retained environments print in this document as Lemma 1, Proposition 1 in order of appearance, whereas the article prints the same items with the numbering of its own full document.  The complete native source of the article as distributed in the arXiv e-print for this exact version is bundled unchanged as \texttt{sources/180972/source0.tex}.
\begin{equation}\label{eq: two short exact sequences}
\begin{gathered}
    0\to \Omega_X\otimes\OO_X(1) \longrightarrow \PP^1(f) \longrightarrow \OO_X(1) \to 0\\
    0\to \mathrm{Sym}^k\Omega_X\otimes\OO_X(1) \longrightarrow \PP^k(f) \longrightarrow \PP^{k-1}(f) \to 0\quad\forall\,k\ge 2
\end{gathered}
\end{equation}

\begin{equation}\label{lem: chern classes tensor product with line bundle}
c_i(\FF\otimes\LL) = \sum_{j=0}^i\binom{r-j}{i-j}c_j(\FF)c_1(\LL)^{i-j}\,.
\end{equation}

\begin{equation}\label{eq: recursive formula}
\begin{cases}
    c_i(\PP^1(f)) = \sum_{j=0}^i\binom{m+1-j}{i-j}(-1)^jc_j\cdot L^{i-j}\\
    c_i(\PP^k(f)) = \sum_{\ell=0}^i\left[\sum_{j=0}^\ell\binom{\binom{m+k-1}{k}-j}{\ell-j}c_j(\mathrm{Sym}^k\Omega_X)\cdot L^{\ell-j}\right]c_{i-\ell}(\PP^{k-1}(f))\,.
\end{cases}
\end{equation}

\begin{lemma}\label{lem: chern classes symmetric power vector bundle small rank}
Let $\FF$ be a vector bundle on a nonsingular surface $X$. Then for any $k\ge 1$,
\[
c(\mathrm{Sym}^k\FF) =
\begin{cases}
    1+k\,c_1(\FF) & \text{if $\rank\FF=1$}\\[3pt]
    1+\binom{k+1}{2}c_1(\FF)+\binom{k+1}{3}\frac{3k+2}{4}c_1^2(\FF)+\binom{k+2}{3}c_2(\FF) & \text{if $\rank\FF=2$}\\[3pt]
    1+\binom{k+2}{3}c_1(\FF)+\frac{1}{2}\binom{k+3}{3}\binom{k+1}{3}c_1^2(\FF)+\binom{k+3}{4}c_2(\FF) & \\[2pt]
    +\binom{k+3}{5}\frac{5\,k^4+20\,k^3-5\,k^2-50\,k-12}{54}c_1^3(\FF) & \\[2pt]
    +\binom{k+3}{5}\frac{5\,k^2+20\,k+6}{6}c_1(\FF)c_2(\FF)+\binom{k+3}{4}\frac{2\,k+3}{5}c_3(\FF) & \text{if $\rank\FF=3$.}
\end{cases}
\]
\end{lemma}

\begin{equation}\label{eq: identity total chern class J_k curve}
c(\PP^k(f)) = 1+(k+1)L-\binom{k+1}{2}c_1\,.
\end{equation}

\begin{proposition}\label{prop: gEDD k-regular surface}
Consider a $k$-regular embedding $f\colon S\hookrightarrow\P^n$ of a nonsingular projective surface $S$. Write $c(S)=1+c_1+c_2$ and $c(\OO_S(1))=1+L$. Then
\begin{equation}\label{eq: gEDD k-regular surface}
\gDD_k(f) = \int_X\alpha_1\,L^2+\alpha_2\,c_1L+\alpha_3\,c_1^2+\alpha_4\,c_2\,,
\end{equation}
where
\begin{align}\label{eq: coefficients alpha}
\begin{split}
    \alpha_1 &= 1+\binom{k+2}{2}+3\binom{k+3}{4}\,,\quad\alpha_2 = -\binom{k+2}{2}\binom{k+2}{3}\,,\\
    \alpha_3 &= \frac{1}{2}\binom{k+1}{3}\binom{k+3}{3}\,,\quad\alpha_4 = \binom{k+3}{4}\,.
\end{split}
\end{align}
\end{proposition}

\begin{proof}
We claim that, for all $k\ge 1$,
\begin{equation}\label{eq: identity total chern class J_k surface}
\resizebox{.9\textwidth}{!}{$\begin{aligned}c_1(\PP^k(f)) &= \binom{k+2}{2}L-\binom{k+2}{3}c_1\\
c_2(\PP^k(f)) &= \frac{1}{2}\binom{k+1}{2}\binom{k+3}{2}L^2-2k\binom{k+3}{4}c_1L+\frac{1}{2}\binom{k+1}{3}\binom{k+3}{3}c_1^2+\binom{k+3}{4}c_2\,.\end{aligned}$}
\end{equation}
Assuming the claim is true, then
\[
\resizebox{\textwidth}{!}{$\begin{aligned}
\gDD_k(f) &= \mu_{k,0}(f)+\mu_{k,1}(f)+\mu_{k,2}(f) = \int_X L^2+c_1(\PP^k(f))L+c_2(\PP^k(f))\\
&= \int_X \left[1+\binom{k+2}{2}+3\binom{k+3}{4}\right]L^2 -\left[\binom{k+2}{3}+2k\binom{k+3}{4}\right]c_1L+\frac{1}{2}\binom{k+1}{3}\binom{k+3}{3}c_1^2+\binom{k+3}{4}c_2\,,
\end{aligned}$}
\]
which, after simplifying, is equal to the desired formula.
To prove the claim, consider the recursive formula \eqref{eq: recursive formula}. The first equation yields the identities $c_1(\PP^1(f))=3L-c_1$ and $c_2(\PP^1(f))=3L^2-2c_1L+c_2$, thus proving \eqref{eq: identity total chern class J_k surface} for $k=1$. Now assume \eqref{eq: identity total chern class J_k surface} true for $k-1$. Observe that $\mathrm{Sym}^k\Omega_S$ has rank $k+1$ for all $k\ge 1$.
Using Lemma \ref{lem: chern classes symmetric power vector bundle small rank}, we get the Chern classes
\begin{equation}
    c_1(\mathrm{Sym}^k\Omega_S) = -\binom{k+1}{2}c_1\,,\quad c_2(\mathrm{Sym}^k\Omega_S) = \frac{3k+2}{4}\binom{k+1}{3}c_1^2+\binom{k+2}{3}c_2\,.
\end{equation}
Applying Lemma \ref{lem: chern classes tensor product with line bundle}, we obtain that
\[
\resizebox{\textwidth}{!}{$\begin{aligned}
    c(\mathrm{Sym}^k\Omega_S\otimes\OO_S(1)) &= \sum_{i=0}^{k+1}c_i(\mathrm{Sym}^k\Omega_S\otimes\OO_S(1)) = \sum_{i=0}^{k+1}\sum_{j=0}^i\binom{k+1-j}{i-j}c_j(\mathrm{Sym}^k\Omega_S)\cdot L^{i-j}\\
    &= 1+(k+1)L+c_1(\mathrm{Sym}^k\Omega_S)+\binom{k+1}{2}L^2+kc_1(\mathrm{Sym}^k\Omega_S)L+c_2(\mathrm{Sym}^k\Omega_S)\\
    &= 1+(k+1)L-\binom{k+1}{2}c_1+\binom{k+1}{2}L^2-k\binom{k+1}{2}c_1L+\frac{3k+2}{4}\binom{k+1}{3}c_1^2+\binom{k+2}{3}c_2\,.
\end{aligned}$}
\]
Using the second exact sequence in \eqref{eq: two short exact sequences} and the induction step, we obtain that
\[
\resizebox{\textwidth}{!}{$\begin{aligned}
    c(\PP^k(f)) &= c(\mathrm{Sym}^k\Omega_S\otimes\OO_S(1))\cdot c(\PP^{k-1}(f))\\
    &= \left(1+(k+1)L-\binom{k+1}{2}c_1+\binom{k+1}{2}L^2-k\binom{k+1}{2}c_1L+\frac{3k+2}{4}\binom{k+1}{3}c_1^2+\binom{k+2}{3}c_2\right)\\
    &\quad\cdot\left(1+\binom{k+1}{2}L-\binom{k+1}{3}c_1+\frac{1}{2}\binom{k}{2}\binom{k+2}{2}L^2-2(k-1)\binom{k+2}{4}c_1L+\frac{1}{2}\binom{k}{3}\binom{k+2}{3}c_1^2+\binom{k+2}{4}c_2\right)\,.
\end{aligned}$}
\]
Expanding the product, one verifies that
\[
\resizebox{\textwidth}{!}{$\begin{aligned}
    c_1(\PP^k(f)) &= \left[(k+1)+\binom{k+1}{2}\right]L-\left[\binom{k+1}{2}+\binom{k+1}{3}\right]c_1 = \binom{k+2}{2}L-\binom{k+2}{3}c_1\,,\\
    c_2(\PP^k(f)) &= \left[(k+2)\binom{k+1}{2}+\frac{1}{2}\binom{k}{2}\binom{k+2}{2}\right]L^2\\
    &\quad-\left[k\binom{k+1}{2}+\binom{k+1}{2}^2+(k+1)\binom{k+1}{3}+2(k-1)\binom{k+2}{4}\right]c_1L\\
    &\quad+\left[\frac{3k+2}{4}\binom{k+1}{3}+\binom{k+1}{2}\binom{k+1}{3}+\frac{1}{2}\binom{k}{3}\binom{k+2}{3}\right]c_1^2+\left[\binom{k+2}{3}+\binom{k+2}{4}\right]c_2\,,
\end{aligned}$}
\]
which, after simplifying, correspond to the identities in \eqref{eq: identity total chern class J_k curve} at the step $k$.
\end{proof}

\end{document}
